\chapter{Copper Deficiency and Excess}
\index{Copper Deficiency and Excess}
\label{sec:Copper Disorders}

\section{Overview}
Copper deficiency is a rare mineral deficiency disorder that occurs after gastric bypass surgery, in malabsorption, prolonged parenteral nutrition, excessive zinc intake (zinc induces metallothionein which blocks intestinal copper absorption), and in Menkes disease. This metabolic disorder involves dysregulation of normal bodily processes, often requiring ongoing monitoring and management. It is increasingly common due to lifestyle and dietary factors. Metabolic disorders involve disruption of normal biochemical processes essential for health. These conditions may result from enzyme deficiencies, hormonal imbalances, or lifestyle factors. Many metabolic disorders are chronic and require ongoing management. Lifestyle modifications are often the cornerstone of treatment.
\section{Causes}
This condition happens because of deficiency of copper, an essential trace element serving as a cofactor for cytochrome c oxidase, superoxide dismutase, lysyl oxidase, dopamine $\beta$-hydroxylase, and ceruloplasmin. Metabolic disorders arise from dysregulation of normal biochemical pathways. This may result from enzyme deficiencies, hormonal imbalances, or lifestyle factors that overwhelm normal homeostatic mechanisms.   
\section{Risk Factors}
\begin{itemize}[noitemsep]
  \item Genetic predisposition and family history
  \item Advancing age
  \item Lifestyle factors (diet, exercise, smoking, alcohol)
  \item Environmental exposures
  \item Underlying medical conditions
  \item Medication use
\end{itemize}
\section{Signs and Symptoms}
\subsection{Common symptoms}
\begin{itemize}[noitemsep]
  \item You may experience microcytic hypochromic anemia (refractory to iron therapy), neutropenia, myeloneuropathy (spastic gait, sensory loss of coordination, peripheral neuropathy mimicking B12 deficiency), and pancytopenia.
  \item Pain or discomfort in the affected area
  \end{itemize}
\subsection{Emergency warning signs}
\begin{itemize}[noitemsep]
  \item Severe or worsening symptoms of copper deficiency and excess require immediate medical evaluation
\end{itemize}
\section{Progression and Stages}
The condition may progress through different stages of severity over time. Early stages often involve mild or subtle symptoms that may be overlooked. Moderate stages involve more noticeable symptoms affecting daily function. Advanced stages may involve significant impairment and complications.
\section{Complications}
\begin{itemize}[noitemsep]
  \item Disease progression leading to worsening symptoms
  \item Secondary infections or injuries
  \item Side effects of treatment
  \item Reduced quality of life
  \item Emotional and psychological impact
\end{itemize}
\section{Diagnosis}
Doctors diagnose this using low serum copper and low ceruloplasmin.

\begin{itemize}[noitemsep]
  \item Comprehensive medical history and physical examination
  \item Laboratory tests to assess organ function
  \item Imaging studies to visualize affected structures
  \item Specialized testing depending on the affected system
  \item Consultation with relevant specialists
\end{itemize}
\section{Prevention}
\begin{itemize}[noitemsep]
  \item Maintain a healthy lifestyle with balanced diet and exercise
  \item Avoid known risk factors
  \item Attend regular medical check-ups for early detection
  \item Follow recommended screening guidelines
\end{itemize}
\section{Treatment Options}
\begin{itemize}[noitemsep]
  \item Treatment focuses on oral copper 2--5 mg/kg/day as copper sulfate. Menkes disease (kinky hair disease) is an X-linked recessive disorder of ATP7A causing impaired intestinal copper absorption, presenting in infancy with progressive neurodegeneration, seizures, hypotonia, failure to thrive, characteristic brittle, depigmented, kinky hair (pili torti), and tortuous cerebral arteries
  \item Subcutaneous copper-histidine may improve outcomes if started early, but Unfortunately, the outlook is often poor. Wilson disease (copper excess) is covered in Chapter~8.
  \item Medications to manage symptoms and address underlying mechanisms
  \item Lifestyle modifications and self-care measures
  \item Physical therapy and rehabilitation
  \item Surgical intervention when indicated
  \item Regular monitoring and follow-up care
\end{itemize}
\section{Living with It}
\begin{itemize}[noitemsep]
  \item Follow your treatment plan as prescribed
  \item Maintain regular follow-up with your healthcare provider
  \item Make necessary lifestyle adjustments
  \item Seek support from family, friends, or support groups
  \item Stay informed about your condition
\end{itemize}
